\chapter{Podsumowanie i wnioski}
\label{chap:wnioski}

W tej pracy sprawdziłem, jak wybrane metody uczenia przez wzmacnianie
radzą sobie z podejmowaniem decyzji w Ultimate Texas Hold'em oraz jak
na jakość uzyskiwanej strategii wpływa sposób reprezentacji stanu
i konstrukcja funkcji nagrody.

W tym celu przygotowałem kompletne środowisko symulacyjne gry,
mechanizm prowadzenia treningu i ewaluacji, agentów bazowych oraz
agentów wykorzystujących DQN i REINFORCE. Następnie przeprowadziłem
eksperyment obejmujący osiem podstawowych konfiguracji, z których każdą
wytrenowałem dla pięciu seedów. Końcową ocenę modeli wykonałem na
wspólnym harmonogramie 100\,000 rozdań.

\section{Realizacja celu pracy}
\label{sec:realizacja-celu-pracy}

Pierwszym celem pracy było przygotowanie środowiska pozwalającego na
prowadzenie powtarzalnych eksperymentów z wykorzystaniem Ultimate
Texas Hold'em. W ramach realizacji tego celu zaimplementowałem przebieg
pełnego rozdania, ocenę układów pokerowych, zasady wykonywania zakładów
oraz ich końcowe rozliczenie.

Na bazie środowiska przygotowałem wspólną infrastrukturę dla agentów.
Pozwala ona uruchamiać te same rozdania dla różnych strategii,
zapisywać wyniki i wykonywane akcje oraz porównywać agentów na
wspólnych harmonogramach talii. Jako punkty odniesienia wykorzystałem
agenta losowego i agenta regułowego.

Kolejnym etapem była implementacja dwóch metod uczenia przez
wzmacnianie. DQN reprezentuje podejście oparte na aproksymacji funkcji
wartości akcji, natomiast REINFORCE bezpośrednio optymalizuje politykę.
Dla obu metod wykorzystałem te same dwa warianty reprezentacji stanu
oraz te same dwa warianty funkcji nagrody, dzięki czemu mogłem
porównywać wpływ poszczególnych elementów konfiguracji.

Za istotną część praktycznego celu pracy uważam również przygotowanie
powtarzalnego procesu eksperymentalnego. Oddzieliłem dane
wykorzystywane podczas treningu, okresowej walidacji oraz końcowej
ewaluacji. Wszystkie modele w końcowym porównaniu oceniłem na tym
samym harmonogramie rozdań. Pozwoliło mi to ograniczyć wpływ
przypadkowej kolejności kart na porównania pomiędzy strategiami.

W rezultacie powstał kompletny proces obejmujący implementację
środowiska, trening agentów, zapis modeli, okresową walidację,
końcową ewaluację oraz analizę wyników.


\section{Wnioski z przeprowadzonych eksperymentów}
\label{sec:wnioski-eksperymenty}

Najbardziej jednoznaczny wynik uzyskałem dla wpływu reprezentacji
stanu na DQN. W obu wariantach funkcji nagrody reprezentacja FEATURES
prowadziła do wyższego końcowego EV niż RAW. Oznacza to, że przy
zastosowanym budżecie treningowym jawne przekazanie agentowi
dodatkowych cech opisujących sytuację pokerową ułatwiało DQN uczenie
lepszej strategii.

Dla REINFORCE nie zaobserwowałem równie wyraźnego efektu reprezentacji
stanu. Przy zysku netto wyniki RAW i FEATURES były zbliżone, natomiast
dla nagrody skalowanej konfiguracja FEATURES osiągnęła wyższy wynik
punktowy, ale jednocześnie charakteryzowała się dużą zmiennością
pomiędzy seedami. Wpływ reprezentacji na REINFORCE okazał się więc
znacznie mniej powtarzalny niż dla DQN.

Drugim analizowanym elementem była funkcja nagrody. Skalowanie wyniku
netto poprawiło rezultat DQN korzystającego z reprezentacji RAW,
jednak przy reprezentacji FEATURES różnica pomiędzy funkcjami nagrody
była niewielka. Dla REINFORCE również nie uzyskałem wyników
pozwalających stwierdzić, że jedna funkcja nagrody jest zawsze lepsza
od drugiej. Skalowanie zmieniało przebieg optymalizacji i skalę
gradientów, ale nie dawało uniwersalnej poprawy jakości końcowej
polityki.

Porównanie samych algorytmów również nie daje jednej odpowiedzi na
pytanie o lepszą metodę. Przy podstawowym budżecie 50\,000 epizodów
REINFORCE osiągał wyższe średnie EV niż DQN dla odpowiadających sobie
konfiguracji. Jednocześnie jego trening był znacznie bardziej zależny
od seedu i w wielu przebiegach prowadził do przedwczesnej zbieżności
decyzji do bardzo prostej polityki.

DQN osiągał niższe średnie wyniki, ale jego zachowanie było bardziej
powtarzalne pomiędzy niezależnymi treningami. Różnica ta była również
widoczna po zwiększeniu budżetu treningowego do 100\,000 epizodów.
Wszystkie pięć seedów wybranej konfiguracji DQN poprawiło końcowy
wynik, podczas gdy dla REINFORCE zmiana była znacznie mniej
regularna.

Ważnym punktem odniesienia pozostają agenci bazowi. Większość
konfiguracji uczenia przez wzmacnianie osiągnęła wynik lepszy od
agenta losowego, co pokazuje, że modele były w stanie nauczyć się
zależności pomiędzy obserwowanym stanem gry i podejmowanymi akcjami.
Żadna z badanych konfiguracji nie zbliżyła się jednak do wyniku
agenta RuleBased, którego EV w końcowej ewaluacji wyniosło
$-0{,}0106$.

Odpowiadając na postawiony we wstępie problem badawczy, wybór
algorytmu, sposobu reprezentacji stanu oraz funkcji nagrody faktycznie
wpływał zarówno na przebieg uczenia, jak i jakość końcowej strategii.
Wpływ tych elementów nie był jednak jednakowy dla obu algorytmów.
Najbardziej powtarzalnym efektem była poprawa DQN po zastosowaniu
reprezentacji FEATURES, natomiast największym problemem REINFORCE
okazała się stabilność procesu uczenia.

Uzyskane wyniki pokazują również, że ocena metody wyłącznie na
podstawie średniego EV może być niewystarczająca. REINFORCE uzyskiwał
wyższe wyniki średnie, ale DQN dawał bardziej powtarzalny rezultat.
W przypadku zastosowania uczenia przez wzmacnianie w problemie
o dużej losowości ważna jest więc nie tylko jakość najlepszej
uzyskanej polityki, ale również to, jak często podobną politykę można
otrzymać po ponownym przeprowadzeniu treningu.